% ============================================================
% STYLE DES SECTIONS
% ============================================================
\titleformat{\section}
  {\large\bfseries\color{maincolor}}
  {}{0em}{}[\color{maincolor}\titlerule]
\titlespacing*{\section}{0pt}{4pt}{3pt}

\titleformat{\subsection}
  {\normalsize\bfseries\color{secondcolor}}
  {}{0em}{}
\titlespacing*{\subsection}{0pt}{5pt}{4pt}

% ============================================================
% COMMANDES CV
% ============================================================

% Puce des listes du corps du CV.
% Passer par une commande plutôt que d'écrire le label en dur dans chaque
% \begin{itemize} permet à une variante de la redéfinir globalement : la
% variante une colonne remplace l'icône FontAwesome par un vrai caractère de
% puce, que les extracteurs de texte (pdftotext, ATS) savent replacer dans le
% flux — l'icône, elle, est rejetée en fin de bloc et vient se coller au
% texte de la puce suivante.
\newcommand{\cvbullet}{\textcolor{maincolor}{\tiny\faIcon{circle}}}

% Sous-puce de détail (deuxième niveau des listes d'expérience). Même logique
% que \cvbullet : passer par une commande laisse une variante ajuster la
% taille globalement — la variante une colonne les compose plus petit pour
% tenir sur une page.
\newcommand{\cvsub}[1]{{\small\color{detailcolor}\itshape #1}}

% Espace laissé sous chaque entrée : c'est lui qui sépare visuellement deux
% expériences (ou deux projets) consécutifs.
\newlength{\cvitemsep}
\setlength{\cvitemsep}{5pt}

% Expérience détaillée : \cvitem{Poste}{Dates}{Lieu}{Contenu}
\newcommand{\cvitem}[4]{
  {\large\textbf{\color{secondcolor}#1}} \hfill {\small\textit{\color{lightgray}#2}}\\[-2pt]
  {\small\textit{\color{accentcolor}#3}}
  \vspace{0.5pt}
  #4
  \par\vspace{\cvitemsep}
}

% Ligne « Environnement » : \cvenv{Techno, Techno, ...}
% Barre verticale à gauche dont la hauteur s'adapte au texte (1 ou N lignes).
\newsavebox{\cvenvbox}
\newcommand{\cvenv}[1]{%
  \par\vspace{0.5pt}\noindent
  \savebox{\cvenvbox}{%
    \parbox{\dimexpr\linewidth-8pt\relax}{%
      \small\textbf{\color{secondcolor}Environnement :} \textbf{\color{detailcolor}#1}%
    }%
  }%
  \textcolor{maincolor}{\rule[-\dp\cvenvbox]{1.6pt}{\dimexpr\ht\cvenvbox+\dp\cvenvbox\relax}}%
  \hspace{6pt}\usebox{\cvenvbox}\par
}

% Projet personnel : \cvproject{Nom}{Stack}
% Contrairement à \cvitem, le nom et la stack partagent une seule ligne : un
% projet n'a ni dates ni lieu à afficher, la ligne « lieu » de \cvitem restait
% donc vide et coûtait une ligne par projet.
\newcommand{\cvproject}[2]{%
  {\large\textbf{\color{secondcolor}#1}}\hfill{\small\textit{\color{accentcolor}#2}}%
  \par\vspace{0.5pt}%
}

% Ligne courte : \cvitemshort{Titre}{Institution | Année}
\newcommand{\cvitemshort}[2]{
  \textbf{\color{secondcolor}#1} \hfill {\small\textit{\color{lightgray}#2}}
  \vspace{1.5pt}
}

% Compétence avec icône
\newcommand{\skill}[1]{
  \textcolor{maincolor}{\faIcon{check-circle}} \textbf{#1}
}

% ============================================================
% LIENS
% ============================================================
\hypersetup{
    colorlinks=true,
    linkcolor=maincolor,
    urlcolor=maincolor,
}
